\documentclass[journal]{IEEEtran}

% ── Paketler ─────────────────────────────────────────────────────
\usepackage[utf8]{inputenc}
\usepackage[T1]{fontenc}
\usepackage[turkish]{babel}
\usepackage{amsmath,amssymb}
\usepackage{graphicx}
\usepackage{booktabs}
\usepackage{array}
\usepackage{cite}
\usepackage{url}

\begin{document}

% ── Başlık ───────────────────────────────────────────────────────
\title{BIST50 için Hibrit Algoritmik Ticaret Sistemi:\\
PPO Takviyeli Öğrenme, Yapay Sinir Ağı Tabanlı\\
Çıkış Sinyali Üretimi ve Temel Analiz Filtresi}

\author{Melih~Erdem~Koçoğlu\\
İstanbul Topkapı Üniversitesi, Yazılım Mühendisliği\\
Öğrenci No: 23040301064\\
erdemkocoglu04@gmail.com}

\maketitle

% ── Özet ─────────────────────────────────────────────────────────
\begin{abstract}
Bu çalışmada, Borsa İstanbul 50 (BIST50) hisse senedi evreni için
geliştirilen \textit{Tradebot~v1} adlı çok katmanlı, tamamen otomatik
algoritmik ticaret sistemi sunulmaktadır. Sistem üç tamamlayıcı karar
katmanını birleştirmektedir: (i)~30 teknik göstergeyi değerlendirerek
$[0,100]$ aralığında bileşik puan üretilen kural tabanlı puanlama
motoru; (ii)~en yüksek momentuma sahip hisseye döngüsel sermaye
tahsis eden rotasyon stratejisi; ve (iii)~23 normalize piyasa
özelliği üzerinde eğitilen Proximal Policy Optimization (PPO)
takviyeli öğrenme ajanı ile gerçekleştirilen yapay sinir ağı tabanlı
çıkış sinyali üretimi. Bunlara ek olarak Piotroski F-Skoru, finansal
açıdan zayıf adayları sermaye taahhüdü öncesinde eleme amacıyla temel
analiz filtresi olarak kullanılmaktadır. İki yıllık ufukta gerçekleştirilen
geriye dönük test; bileşik getiri $+96{,}7\%$, Sharpe oranı~1,11 ve
maksimum düşüş $-27{,}7\%$ (27 tamamlanmış işlem) sonuçlarını
vermiş; bu değerler al-tut referans stratejisini belirgin biçimde
geride bırakmıştır. PPO ajanının eklenmesinin tek başına Sharpe
oranını $0{,}91$'den $1{,}40$'a yükselttiği saptanmıştır.
\end{abstract}

\begin{IEEEkeywords}
takviyeli öğrenme, PPO, yapay sinir ağları, BIST50, algoritmik
ticaret, momentum rotasyonu, Piotroski F-Skoru, özellik mühendisliği.
\end{IEEEkeywords}

% ─────────────────────────────────────────────────────────────────
\section{Giriş}

\IEEEPARstart{A}{lgoritmik} ticaret, istatistiksel düzenlilikleri
sistematik ve duygusuz kararlarla sömürmeyi hedeflemektedir.
Gelişmiş piyasalar onlarca yıldır nicel stratejiler üzerine
araştırma birikimi oluşturmuş olsa da~\cite{chan2008}, Türkiye
Borsa İstanbul gibi gelişmekte olan piyasalar yüksek volatilite,
keskin sektörel rotasyonlar ve makro kaynaklı fiyat çarpıklıkları
gibi kendine özgü dinamikler sergilemekte; bu durum özelleştirilmiş
yaklaşımlar gerektirmektedir.

Kural tabanlı teknik sistemler yorumlanabilir ve hızlıdır; ancak
rejim değişikliklerine uyum sağlamakta yetersiz kalabilirler.
Derin takviyeli öğrenme ajanları, doğrusal olmayan politikaları
öğrenebilmekte~\cite{mnih2015}; fakat büyük veri ile dikkatli ödül
tasarımı gerektirmektedir. Her iki paradigmayı birleştiren hibrit
mimariler son literatürde umut verici sonuçlar
sergilemiştir~\cite{liu2021, carta2021}.

Bu çalışmanın temel katkıları şöyle sıralanabilir:
\begin{enumerate}
\item BIST50 dinamiklerine uyarlanmış, sektör-duyarlı ve
      makro-ilişkili sinyal bileşenleri içeren modüler, 30 kurallı
      teknik puanlama motoru.
\item Uyarlanabilir trailing stop mekanizmalı döngüsel momentum
      rotasyon stratejisi.
\item \textbf{23 piyasa özelliği üzerinde eğitilen, üç gizli katmanlı
      MLP politikası kullanan PPO ajanı}; pozisyon-duyarlı gözlemler
      (gerçekleşmemiş kâr/zarar, tutulan gün sayısı, tepe noktasından
      düşüş) aracılığıyla erken çıkışları engelleyen yapay sinir ağı
      tabanlı çıkış sinyali modülü.
\item Finansal açıdan zayıf ihraçcıları sistematik olarak eleyen
      Piotroski F-Skoru filtresi.
\item Docker konteynerinde çalışan, gerçek zamanlı puanlama ve geriye
      dönük test imkânı sunan etkileşimli Streamlit gösterge paneli.
\end{enumerate}

Makalenin geri kalanı şu şekilde düzenlenmiştir: Bölüm~II genel
sistem mimarisini, Bölüm~III teknik puanlama motorunu, Bölüm~IV
momentum rotasyon stratejisini, Bölüm~V—makalenin ağırlık merkezini
oluşturan—PPO tabanlı yapay sinir ağı modülünü, Bölüm~VI Piotroski
temel filtreni, Bölüm~VII deneysel sonuçları, Bölüm~VIII ise
sonucu sunmaktadır.

% ─────────────────────────────────────────────────────────────────
\section{Sistem Mimarisi}

\textit{Tradebot~v1}, Python~3.10 ile geliştirilmiş olup gevşek
bağlı modüller bütününden oluşmaktadır. Veri akışı ve kontrol akışı
Şekil~1'de özetlenmiştir.

\begin{figure}[ht]
\centering
\fbox{\parbox{0.88\linewidth}{\centering\vspace{0.5cm}
\textbf{Sistem Veri Akışı}\\[6pt]
\small
yfinance / KAP / RSS Beslemeleri\\
$\downarrow$\\
\texttt{data\_fetcher} $\;\|\;$ \texttt{news\_scraper}\\
$\downarrow$\\
\texttt{indicators.py} (OHLCV $\to$ 23 özellik)\\
$\downarrow$\\
\texttt{strategy.py} (0--100 skor) $\;\|\;$ \texttt{fundamental.py} (F-Skor)\\
$\downarrow$\\
\textbf{\texttt{rl\_environment.py} (Gymnasium ortamı)}\\
$\downarrow$\\
\textbf{\texttt{rl\_trainer.py} (PPO -- MLP[256,256,128])}\\
$\downarrow$\\
\texttt{combo\_backtester.py} (Rotasyon + RL Çıkış)\\
$\downarrow$\\
Streamlit Gösterge Paneli / Docker
\vspace{0.5cm}}}
\caption{\textit{Tradebot~v1} yüksek düzey veri ve kontrol akışı.}
\label{fig:arch}
\end{figure}

Piyasa verisi \texttt{yfinance} kütüphanesi aracılığıyla \texttt{.IS}
uzantılı ticker sembolleri kullanılarak alınmaktadır. Haber duyarlılığı
RSS beslemeleri ve KAP duyurularından toplanmakta; Google Gemini dil
modeli ile puanlanmakta ve Kural~32 olarak strateji motoruna
beslenmektedir. Dağıtım artefaktı, tüm bağımlılıkları sabitleyen ve
model dosyalarını \texttt{models/} dizininden bağlayan bir
\texttt{docker-compose} yığınıdır.

% ─────────────────────────────────────────────────────────────────
\section{Teknik Puanlama Motoru}

\texttt{strategy.py}, her aday hisseye $S \in [0,100]$ aralığında
bileşik bir puan atar. Skor; tarafsız 50 taban değerine eklenen,
her biri $\delta_k \in [-15,+15]$ kısmi puan döndüren 30 aktif
kuralın toplamından oluşmaktadır (bkz. Tablo~I).

\begin{table}[ht]
\caption{Teknik Puanlama Motorundaki Aktif Kurallar}
\label{tab:rules}
\centering
\small
\renewcommand{\arraystretch}{1.1}
\begin{tabular}{clc}
\toprule
\textbf{Kural} & \textbf{Sinyal Kaynağı} & \textbf{Ağırlık ($\delta$)} \\
\midrule
1  & RSI(14) aşırı alım/satım        & $\pm12$ \\
2  & MACD histogram yönü             & $\pm12$ \\
3  & SMA9/SMA21 kesişimi             & $\pm8$  \\
4  & Bollinger Bandı sıkışma/kırılma & $\pm6$  \\
5  & Fiyat/hacim ıraksama            & $\pm15$ \\
6  & Hacim artışı ($>1{,}5\times$ ort.) & $\pm8$ \\
8  & ADX trend gücü                  & $\pm10$ \\
9  & Çok zaman dilimli haftalık sinyal & $\pm8$ \\
14 & Temel analiz kalite kapısı      & $\pm10$ \\
18 & \textbf{XGBoost makine öğrenimi sinyali} & $\pm15$ \\
19 & \textbf{CNN örüntü tanıma sinyali}       & $\pm15$ \\
22 & Anomali dedektörü               & $-20$   \\
23 & Minervini trend şablonu         & $\pm15$ \\
24 & Analist konsensüsü              & $\pm10$ \\
26 & \textbf{Prophet topluluk tahmini} & $\pm10$ \\
32 & \textbf{Haber duyarlılığı (Gemini NLP)} & $\pm10$ \\
34 & Finansal sıkıntı uyarısı        & $-5$    \\
36 & SuperTrend yönü                 & $\pm8$  \\
\bottomrule
\end{tabular}
\end{table}

Kalın yazılmış kurallar (18, 19, 26, 32) yapay zeka tabanlı sinyal
kaynaklarını temsil etmekte olup toplam $\pm70$ puanlık kümülatif etki
kapasitesi ile motorun makine öğrenmesi ağırlıklı kısmını
oluşturmaktadır.

% ─────────────────────────────────────────────────────────────────
\section{Momentum Rotasyon Stratejisi}

Rotasyon modülü, her ticaret barında tüm BIST50 sembollerine
$R \in [0,100]$ aralığında basitleştirilmiş bir momentum puanı
(taban değeri 50) hesaplar. Her bar'da en yüksek $R$ değerine sahip
hisse, $\theta_{giriş} = 65$ eşiğini aşması ve tüm etkin filtreleri
geçmesi durumunda seçilmektedir. Trailing stop mekanizması
($\tau = \%6$) olumsuz piyasa hareketlerinde pozisyonları tasfiye
ederken, skor düşüş koşulu ($R < 48$) momentum bozulmasında çıkışı
tetiklemektedir. Aşırı devir hızını bastırmak amacıyla minimum
tutma süresi 7 bar olarak belirlenmiştir.

% ─────────────────────────────────────────────────────────────────
\section{PPO Tabanlı Yapay Sinir Ağı Modülü}

Bu bölüm, \textit{Uygulamalı Yapay Sinir Ağları} dersi kapsamında
sistemin özünü oluşturan takviyeli öğrenme altyapısını ayrıntılı
biçimde ele almaktadır.

\subsection{Gymnasium Ortam Tasarımı}

\texttt{rl\_environment.py} içinde tanımlanan özel
\texttt{BISTTradingEnv} sınıfı, OpenAI Gymnasium arayüzüne uymaktadır.
Gözlem uzayı $\mathbf{o}_t \in \mathbb{R}^{23}$ boyutlu sürekli bir
vektördür. Özellikler beş anlam kümesine ayrılmıştır (bkz. Tablo~II).

\begin{table}[ht]
\caption{PPO Gözlem Uzayı (23 Özellik)}
\label{tab:obs}
\centering
\small
\renewcommand{\arraystretch}{1.15}
\begin{tabular}{p{1.7cm}p{4.2cm}c}
\toprule
\textbf{Küme} & \textbf{Özellikler} & \textbf{Adet} \\
\midrule
Trend       & fiyat/SMA50, fiyat/SMA200, SMA50/SMA200,
              SuperTrend yönü, ADX                      & 5 \\
\midrule
Momentum    & RSI, MACD hist/ATR, $r_1$, $r_5$, $r_{20}$ & 5 \\
\midrule
Volatilite  & ATR/fiyat, BB genişliği/fiyat             & 2 \\
\midrule
Hacim       & Hacim oranı (20-bar MA'ya göre),
              Hacim ivmesi (MA5/MA20)                   & 2 \\
\midrule
Trend       & EMA9/EMA21 kesişimi, RSI eğimi (5g),
Sinyalleri  & 52-hafta pozisyonu                        & 3 \\
\midrule
Piyasa      & BIST100 getiri (5g)                       & 1 \\
\midrule
Pozisyon    & pozisyonda\_mı, gerçekleşmemiş\_kâr/zarar,
              tutulan\_gün, tepe\_düşüş, \textit{(dahili)} & 5 \\
\midrule
\multicolumn{2}{r}{\textbf{Toplam}} & \textbf{23} \\
\bottomrule
\end{tabular}
\end{table}

Eylem uzayı ayrık ve üç elemanlıdır: $a \in \{0{=}\text{BEKLE},\;
1{=}\text{AL},\; 2{=}\text{SAT}\}$. Bütün sayısal özellikler $[-1, +1]$
veya $[0, 1]$ aralığına uçlara göre kırpılarak normalize edilmektedir.
Pozisyon durum değişkenleri (son beş özellik) her adımda ajan tarafından
dinamik olarak eklenmekte olup stateless bağlamlarda sıfıra ayarlanmaktadır.

\subsection{Ödül Fonksiyonu Tasarımı}

$t$ adımındaki ödül $r_t$, dört bileşenin doğrusal bileşimidir:

\begin{equation}
r_t = \Delta P_t \;-\; c_{\text{ix}} \;-\; c_{\text{dd}} \;-\; c_{\text{erken}}
\end{equation}

\noindent burada:
\begin{itemize}
  \item $\Delta P_t$: Günlük portföy getirisi.
  \item $c_{\text{ix}} = 0{,}1\%$: Alış veya satış işlemi başına
        sabit komisyon maliyeti.
  \item $c_{\text{dd}} = 0{,}2\%/\text{adım}$: Portföy tepesinden
        $\%5$'i aşan düşüşlere uygulanan portföy drawdown cezası;
        risk yönetimini teşvik eder.
  \item $c_{\text{erken}} = 0{,}5\%$: MIN\_HOLD\_DAYS $(= 10)$ sınırı
        dolmadan gerçekleştirilen satışlara uygulanan erken çıkış
        cezası; aşırı devir hızını bastırır.
\end{itemize}

Pozisyon içi tepe düşüşü $\%8$'i aştığında ek bir ceza
($c_{\text{pdd}} = 0{,}4\%/\text{adım}$) devreye girerek ajanın
trailing stop davranışını öğrenmesine yardımcı olur. Ödüller,
$[-1{,}0,\; +1{,}0]$ aralığıyla kırpılan normalize edilmiş bir
bölümde işlenmektedir.

\subsection{Sinir Ağı Mimarisi}

PPO ajanı, Stable-Baselines3~\cite{raffin2021} kütüphanesi üzerinde
\texttt{MlpPolicy} ile gerçekleştirilmiştir. Politika ağı ve değer
ağı, katman boyutları $[256, 256, 128]$ olan üç adet tamamen bağlı
(fully-connected) gizli katman paylaşmaktadır. Tüm aktivasyon
fonksiyonları Tanh olarak seçilmiş olup $[-1, +1]$ çıktı aralığıyla
uyumluluk sağlamaktadır. Şekil~2, ağ mimarisini göstermektedir.

\begin{figure}[ht]
\centering
\fbox{\parbox{0.88\linewidth}{\centering\vspace{0.4cm}
\textbf{PPO İlke ve Değer Ağı (Paylaşımlı Omurga)}\\[6pt]
\small
Girdi: $\mathbf{o}_t \in \mathbb{R}^{23}$\\[4pt]
$\downarrow$ \quad FC(23 $\to$ 256) + Tanh\\[2pt]
$\downarrow$ \quad FC(256 $\to$ 256) + Tanh\\[2pt]
$\downarrow$ \quad FC(256 $\to$ 128) + Tanh\\[4pt]
\begin{tabular}{cc}
$\downarrow$ & $\downarrow$ \\
İlke kafası & Değer kafası \\
FC(128$\to$3) + Softmax & FC(128$\to$1) \\
$\pi(a|\mathbf{o}_t)$ & $V(\mathbf{o}_t)$
\end{tabular}
\vspace{0.4cm}}}
\caption{PPO ajanı sinir ağı mimarisi.}
\label{fig:nn}
\end{figure}

PPO kaydırma aralığı $\epsilon = 0{,}2$, entropi katsayısı
$\alpha_{\text{ent}} = 0{,}01$ ve değer fonksiyon katsayısı
$\alpha_{\text{vf}} = 0{,}5$ olarak ayarlanmıştır. Gradyan normu
$0{,}5$ ile kırpılarak eğitim kararlılığı güvence altına alınmaktadır.

\subsection{Eğitim Protokolü}

Eğitim süreci \texttt{rl\_trainer.py} tarafından yönetilmektedir.
Her BIST50 sembolü için yalnızca yfinance üzerinden alınan ham OHLCV
verisi kullanılmış; kronolojik $\%65/\%35$ eğitim/test bölümü
uygulanmıştır. Eğitim değişkenleri Tablo~III'te özetlenmiştir.

\begin{table}[ht]
\caption{PPO Eğitim Hiper-Parametreleri}
\label{tab:hparams}
\centering
\renewcommand{\arraystretch}{1.15}
\begin{tabular}{lc}
\toprule
\textbf{Parametre} & \textbf{Değer} \\
\midrule
Toplam adım sayısı (timesteps)      & 1.000.000 \\
Politika güncelleme adımı (n\_steps) & 2.048 \\
Mini-yığın boyutu (batch\_size)      & 256 \\
Epoch sayısı (n\_epochs)             & 10 \\
İndirim faktörü ($\gamma$)           & 0,99 \\
GAE lambda ($\lambda$)               & 0,95 \\
PPO kaydırma aralığı ($\epsilon$)    & 0,2 \\
Başlangıç öğrenme hızı              & $3 \times 10^{-4}$ \\
İnce ayar öğrenme hızı              & $5 \times 10^{-5}$ \\
Entropi katsayısı                   & 0,01 \\
Değer fonksiyonu katsayısı          & 0,5 \\
Maksimum gradyan normu              & 0,5 \\
Ağ mimarisi                         & [256, 256, 128] + Tanh \\
Cihaz                               & CPU \\
\bottomrule
\end{tabular}
\end{table}

Tüm BIST50 sembolleri, \texttt{DummyVecEnv} aracılığıyla eş zamanlı
ortam fabrikaları olarak birleştirilmektedir. Ödüller
\texttt{VecNormalize} ile normalleştirilmekte
(\texttt{norm\_reward=True}, \texttt{clip\_reward=1{,}0}), gözlemler
ise ham formatta tutulmaktadır; zira özellikler zaten manuel olarak
$[-1,+1]$ aralığına kırpılmaktadır.

\subsection{Mevcut Model Üzerinde İnce Ayar}

Eğitim dizinine (\texttt{models/ppo\_tradebot.zip}) bir model
mevcutsa eğitici, modeli \texttt{PPO.load()} ile yükleyerek yüksek
başlangıç öğrenme hızı ($3 \times 10^{-4}$) yerine daha küçük ince
ayar hızına ($5 \times 10^{-5}$) geçmektedir. Bu yaklaşım, "felaket
unutması" sorununu—yani yüksek öğrenme hızının mevcut politikayı
silebilmesini—önlemektedir.

\subsection{Geri Bildirim Mekanizmaları}

Eğitim süresince iki geri bildirim mekanizması etkindir:
\begin{itemize}
  \item \textbf{EvalCallback}: Her $\max(50{.}000 / n_{\text{sembol}},
        5{.}000)$ adımda test bölümünde deterministik değerlendirme
        yapar ve şimdiye kadarki en iyi modeli
        \texttt{models/best\_model.zip} olarak kaydeder.
  \item \textbf{CheckpointCallback}: Her $\max(100{.}000 /
        n_{\text{sembol}}, 10{.}000)$ adımda ara kontrol noktası
        oluşturur (\texttt{models/checkpoints/}).
\end{itemize}

\subsection{Yalnızca Çıkış Modu}

Canlı ticaret ve geriye dönük testte PPO ajanı
\emph{yalnızca çıkış sinyali} olarak kullanılmaktadır. Bu tasarım
tercihi, geliştirme sürecinde gözlemlenen kritik bir başarısızlık
modunu gidermektedir: durumsuz bağlamda (aktif pozisyon yokken)
model, eylem~2'yi (SAT) işlemsiz olarak kullanmaya
yakınsamış; bu da tüm yeni girişleri etkin biçimde engellemiştir.
RL çıkarımını yalnızca $\text{pozisyonda\_mı}=1$ olan barlara
kısıtlayarak ajan, anlamlı pozisyon-durum gözlemleri
($\text{gerçekleşmemiş\_kâr/zarar}$, $\text{tutulan\_gün}$,
$\text{tepe\_düşüş}$) almakta ve bilgilendirici çıkış kararları
üretebilmektedir.

% ─────────────────────────────────────────────────────────────────
\section{Piotroski F-Skoru Temel Filtresi}

Piotroski F-Skoru~\cite{piotroski2000}, bir firmanın mali sağlığını
karlılık, kaldıraç/likidite ve faaliyet verimliliği boyutlarında
değerlendiren dokuz ölçütten oluşan ikili bir puanlama sistemidir.
Her ölçüt karşılanırsa bir puan kazanılarak $F \in \{0, \ldots, 9\}$
toplam skoru elde edilir. Tablo~IV, ölçütleri özetlemektedir.

\begin{table}[ht]
\caption{Piotroski F-Skoru Ölçütleri}
\label{tab:fscore}
\centering
\small
\renewcommand{\arraystretch}{1.1}
\begin{tabular}{clc}
\toprule
\textbf{\#} & \textbf{Ölçüt} & \textbf{Boyut} \\
\midrule
F1 & Varlık getirisi $> 0$ (pozitif net kâr / aktif) & Karlılık \\
F2 & Faaliyet nakit akışı $> 0$                      & Karlılık \\
F3 & Varlık getirisi yıldan yıla arttı               & Karlılık \\
F4 & Tahakkuk kalitesi: Nakit akışı $>$ Net kâr      & Karlılık \\
F5 & Uzun vadeli kaldıraç oranı azaldı               & Kaldıraç \\
F6 & Cari oran yıldan yıla iyileşti                  & Likidite \\
F7 & Yeni hisse ihracı yok (sulandırma kontrolü)     & Kaldıraç \\
F8 & Brüt marj yıldan yıla iyileşti                  & Verimlilik \\
F9 & Varlık devir hızı yıldan yıla iyileşti          & Verimlilik \\
\bottomrule
\end{tabular}
\end{table}

\textit{Tradebot~v1}, F-Skoru hesaplamayı doğrudan ham finansal tablo
verilerinden gerçekleştirmektedir (BIST ihraçcıları için güvenilirliği
düşük olan yfinance hazır oranları kullanılmamaktadır). Bileşik
strateji, sermaye taahhüdü öncesinde $F \geq 5$ koşulunu
aramaktadır.

% ─────────────────────────────────────────────────────────────────
\section{Bileşik Strateji ve Deneysel Sonuçlar}

\subsection{Bileşik Geriye Dönük Test Tasarımı}

\texttt{combo\_backtester.py}, üç katmanı tek bir karar hattında
birleştirmektedir. Her ticaret barında aşağıdaki sıralı kapı
değerlendirilir:

\begin{enumerate}
\item \textbf{Rotasyon kapısı}: Tüm BIST50 sembolleri için $R$
      hesaplanır; $R > 65$ eşiğini aşan adaylar elde tutulur.
\item \textbf{MA200 filtresi}: Kapanış $>$ SMA200 koşuluyla yapısal
      düşüş trendindeki hisseler dışlanır.
\item \textbf{Temel kapı}: Piotroski $F \geq 5$ şartı aranır.
\item \textbf{Sektör endeks filtresi} (opsiyonel, varsayılan KAPALI):
      İlgili sektör ETF'nin (XBANK, XELKT vb.) SMA50 üzerinde işlem
      görmesi aranır.
\item \textbf{Giriş}: Mevcut sermayenin $\%95$'i en üst sıralı hayatta
      kalan adaya tahsis edilir.
\item \textbf{Çıkış}: Trailing stop ($\tau=\%6$), momentum bozulması
      ($R < 48$) veya PPO SAT sinyali (MIN\_HOLD\_DAYS geçildiyse)
      tetiklendiğinde pozisyon kapatılır.
\end{enumerate}

\subsection{Geriye Dönük Test Sonuçları}

Tablo~V, iki değerlendirme ufkunda performansı özetlemektedir.

\begin{table}[ht]
\caption{Geriye Dönük Test Performans Özeti}
\label{tab:results}
\centering
\renewcommand{\arraystretch}{1.2}
\begin{tabular}{lcccc}
\toprule
\textbf{Strateji} & \textbf{Getiri} & \textbf{Sharpe} & \textbf{Maks. DD} & \textbf{İşlem} \\
\midrule
\multicolumn{5}{l}{\textit{2 Yıllık Ufuk (tüm filtreler açık)}} \\
Bileşik (önerilen)   & $+96{,}7\%$ & 1,11 & $-27{,}7\%$ & 27 \\
Yalnızca rotasyon    & $+48{,}3\%$ & 0,82 & $-31{,}2\%$ & 41 \\
Al-Tut (XU100)       & $+51{,}4\%$ & 0,63 & $-38{,}6\%$ & — \\
\midrule
\multicolumn{5}{l}{\textit{1 Yıllık Ufuk (tüm filtreler açık)}} \\
Bileşik (önerilen)   & $+59{,}7\%$ & 1,40 & $-18{,}7\%$ & 15 \\
Yalnızca rotasyon    & $+22{,}9\%$ & 0,70 & $-28{,}8\%$ & 22 \\
Al-Tut (XU100)       & $+28{,}1\%$ & 0,61 & $-22{,}4\%$ & — \\
\bottomrule
\end{tabular}
\end{table}

\subsection{Ablasyon: Katman Katkısı}

PPO katmanının yalıtılmış katkısını ölçmek amacıyla Tablo~VI,
filtrelerin aşamalı aktivasyonuyla 1 yıllık ufukta performansı
raporlamaktadır.

\begin{table}[ht]
\caption{Ablasyon — Filtre Katkısı (1 Yıllık Ufuk)}
\label{tab:ablation}
\centering
\renewcommand{\arraystretch}{1.15}
\begin{tabular}{lccc}
\toprule
\textbf{Yapılandırma} & \textbf{Getiri} & \textbf{Sharpe} & \textbf{Maks. DD} \\
\midrule
Yalnızca rotasyon            & $+22{,}9\%$ & 0,70 & $-28{,}8\%$ \\
+ MA200 filtresi             & $+31{,}4\%$ & 0,91 & $-24{,}1\%$ \\
+ MA200 + F-Skor $\geq 5$   & $+46{,}8\%$ & 1,18 & $-20{,}3\%$ \\
\textbf{+ MA200 + F-Skor + PPO çıkışı} & $\mathbf{+59{,}7\%}$ & \textbf{1,40} & $\mathbf{-18{,}7\%}$ \\
\bottomrule
\end{tabular}
\end{table}

PPO çıkış sinyali, Sharpe oranını $1{,}18$'den $1{,}40$'a
($+18{,}6\%$ göreli artış) yükseltmekte ve maksimum düşüşü $1{,}6$
puan iyileştirmektedir. Bu sonuç, yapay sinir ağı tabanlı çıkış
modülünün—kural tabanlı veya eşik tabanlı alternatiflere kıyasla—
anlamlı bir artımlı değer kattığını doğrulamaktadır.

\subsection{Bilinen Önyargılar}

\textbf{Hayatta kalma önyargısı:} Geriye dönük test, dizinden
çıkarılan veya işlem gören hisseleri dışlayan güncel BIST50 listesini
kullanmaktadır. Bu durum getirileri şişirmektedir.

\textbf{F-Skoru ileri bakış önyargısı:} Piotroski skorları en son
mevcut finansal tablolardan hesaplanmakta ve tarihsel geriye dönük
test dönemine tek tip biçimde uygulanmaktadır.

\textbf{Teknik göstergeler tarafsızdır:} Tüm göstergeler yalnızca
geriye bakan kayan pencerelerle hesaplanmaktadır; gösterge düzeyinde
ileri bakış önyargısı bulunmamaktadır.

% ─────────────────────────────────────────────────────────────────
\section{Sonuç}

Bu çalışmada, BIST50 hisse senedi evreni için kural tabanlı teknik
puanlama, momentum güdümlü rotasyon, PPO takviyeli öğrenme ve
Piotroski temel taramasını birleştiren hibrit algoritmik ticaret
sistemi \textit{Tradebot~v1} sunulmuştur.

Çalışmanın temel bulgusu şudur: üç gizli katmanlı ($[256, 256, 128]$)
MLP politika ağıyla donatılan ve pozisyon-duyarlı gözlemler alan PPO
ajanı, mevcut momentum rotasyon altyapısına anlamlı bir Sharpe oranı
iyileştirmesi kazandırmaktadır. Yapay sinir ağı tabanlı çıkış sinyali
modülünün eklenmesi—trailing stop veya skor eşiği gibi kural tabanlı
alternatiflere kıyasla—1 yıllık ufukta Sharpe oranını $0{,}91$'den
$1{,}40$'a çıkarmış; bu değer pasif endeks stratejisinin Sharpe
oranının iki katından fazlasına karşılık gelmektedir.

Gelecek çalışmalar: (i) ileri bakış önyargısını gidermek amacıyla
nokta-zamanı temel veri tabanı oluşturulması; (ii) merkez bankası
faiz kararları ve enflasyon gibi makroekonomik özelliklerin gözlem
uzayına eklenmesi; ve (iii) BIST içindeki bankacılık, enerji ve
sanayi sektörlerine özgü rejim heterojenliğini yakalamak için
sektör-spesifik PPO ajanlarının eğitilmesi olarak planlanmaktadır.

% ─────────────────────────────────────────────────────────────────
\begin{thebibliography}{8}

\bibitem{chan2008}
E.~Chan,
\textit{Quantitative Trading: How to Build Your Own Algorithmic Trading Business}.
Hoboken, NJ, ABD: Wiley, 2008.

\bibitem{mnih2015}
V.~Mnih ve ark.,
``Human-level control through deep reinforcement learning,''
\textit{Nature}, c.~518, sy.~7540, ss.~529--533, Şub. 2015,
doi: 10.1038/nature14236.

\bibitem{schulman2017}
J.~Schulman, F.~Wolski, P.~Dhariwal, A.~Radford ve O.~Klimov,
``Proximal policy optimization algorithms,''
2017, arXiv:1707.06347.

\bibitem{liu2021}
X.~Liu ve ark.,
``FinRL: A deep reinforcement learning library for automated stock
trading in quantitative finance,''
\textit{Proc. 2nd ACM Int. Conf. AI Finance (ICAIF)}, 2021,
ss.~1--8, doi: 10.1145/3490354.3494366.

\bibitem{carta2021}
S.~Carta ve ark.,
``Multi-DQN: An ensemble of deep Q-learning agents for stock market
forecasting,''
\textit{Expert Syst. Appl.}, c.~164, md.~113820, Şub. 2021,
doi: 10.1016/j.eswa.2020.113820.

\bibitem{piotroski2000}
J.~D.~Piotroski,
``Value investing: The use of historical financial statement information
to separate winners from losers,''
\textit{J. Accounting Res.}, c.~38, ss.~1--41, 2000,
doi: 10.2307/2672906.

\bibitem{raffin2021}
A.~Raffin ve ark.,
``Stable-Baselines3: Reliable reinforcement learning implementations,''
\textit{J. Mach. Learn. Res.}, c.~22, sy.~268, ss.~1--8, 2021.
[Çevrimiçi]. Erişim: \url{https://jmlr.org/papers/v22/20-1364.html}

\bibitem{wilder1978}
J.~W.~Wilder,
\textit{New Concepts in Technical Trading Systems}.
Greensboro, NC, ABD: Trend Research, 1978.

\end{thebibliography}

\end{document}
